\abstract{
% English abstract.

Handover (HO) management is one of the most important mobility functions in
fifth-generation (5G) and emerging sixth-generation (6G) radio access networks. When handovers are configured or monitored
poorly they cause radio link failures (RLF), ping-pong effects, and a worse
experience for the user, and as networks become denser every user goes through
many more of them, which makes reliable HO monitoring increasingly important.
This thesis studies one question in particular: does HO monitoring become more
reliable when short-term anomalies and long-term concept drift are handled
together instead of separately? The work is scoped to intra-RAT, inter-gNB
handover in 5G Standalone (SA) networks, meaning handover between cells that use
the same radio access technology (RAT) but are served by different base stations
(next-generation NodeBs, gNBs). This is among the most common mobility
situations in dense deployments.

A practical obstacle shapes the methodology. No public dataset exposes the HO
control parameters (time-to-trigger, hysteresis, A3 offset) together with the
mobility key performance indicators (KPIs) they produce, so the main data comes
from a custom 5G New Radio (NR) SA handover simulator built to follow the
relevant 3rd Generation Partnership Project (3GPP) procedures. Its radio layer is calibrated against the real NordicDat drive-test
data, with close agreement for reference signal received power (RSRP) and
signal-to-interference-plus-noise ratio (SINR) and a documented residual
mismatch for reference signal received quality (RSRQ). Real and simulated
samples are never mixed, so the benchmark stays clean.

On this data the thesis builds and compares four layers: unsupervised anomaly
detection on HO KPI windows, concept-drift detection on KPI streams, a
drift-aware monitor that retrains the anomaly detector when drift is detected,
and a configuration-performance surrogate that links HO parameters to mobility
outcomes.

Four findings stand out. Unsupervised detectors catch structured faults well,
with the area under the precision--recall curve (PR-AUC) reaching $0.858$, but
miss most subtle ones. Every injected drift is
eventually detected, so the real design choice is a trade-off between detection
latency and false alarms rather than whether drift is seen at all. Retraining the
monitor naively on recent data backfires: it absorbs post-drift anomalies into
``normal'' and roughly halves detection quality, a self-poisoning effect. A
simple self-supervised filter that drops the most suspicious windows before
retraining repairs this and more than triples drift-phase PR-AUC. Finally, a
gradient-boosting surrogate predicts the handover success rate (HOSR) accurately
($R^2 = 0.96$) and, wrapped
in conformal prediction intervals, drives an end-to-end detect, recommend, and
recover loop with a model-predicted HOSR gain of $+0.62$. The overall message is
that how a monitor is kept up to date matters more than whether it is updated at
all.

}{
% Streszczenie w języku polskim.

Zarządzanie przełączeniem (handover, HO) jest jedną z najważniejszych funkcji
mobilności w sieciach radiowych piątej generacji (5G) oraz powstających
sieciach szóstej generacji (6G). Źle
skonfigurowane lub źle monitorowane przełączenia powodują utratę połączenia
radiowego (radio link failure, RLF), efekt ping-pong i pogorszenie jakości użytkowania, a wraz z
zagęszczaniem sieci każdy użytkownik wykonuje ich znacznie więcej, co czyni
niezawodne monitorowanie HO coraz ważniejszym. Praca bada jedno pytanie: czy
monitorowanie HO staje się bardziej niezawodne, gdy krótkoterminowe anomalie i
długoterminowy dryf pojęć są traktowane łącznie, a nie osobno. Zakres ograniczono
do przełączeń intra-RAT, inter-gNB w sieciach 5G w trybie Standalone (SA), czyli
przełączeń między komórkami korzystającymi z tej samej technologii dostępu
radiowego (radio access technology, RAT), lecz obsługiwanymi przez różne stacje
bazowe (next-generation NodeB, gNB). Jest to jeden z najczęstszych scenariuszy
mobilności w gęstych sieciach.

Metodologię kształtuje praktyczne ograniczenie. Żaden publiczny zbiór danych nie
udostępnia jednocześnie parametrów sterujących HO (time-to-trigger, histereza,
offset A3) oraz wynikających z nich kluczowych wskaźników wydajności (key
performance indicators, KPI) mobilności, dlatego główne dane pochodzą z własnego
symulatora przełączeń 5G New Radio (NR) SA, zbudowanego zgodnie z odpowiednimi
procedurami organizacji 3rd Generation Partnership Project (3GPP). Jego warstwę radiową skalibrowano względem
rzeczywistych danych z testów drogowych NordicDat, z dobrą zgodnością dla mocy
odbieranego sygnału referencyjnego (reference signal received power, RSRP) i
stosunku sygnału do interferencji i szumu (signal-to-interference-plus-noise
ratio, SINR) oraz udokumentowaną resztkową rozbieżnością dla jakości odbieranego
sygnału referencyjnego (reference signal received quality, RSRQ). Danych
rzeczywistych
nigdy nie łączono z symulowanymi, dzięki czemu benchmark pozostaje czysty.

Na tych danych zbudowano i porównano cztery warstwy: nienadzorowaną detekcję
anomalii w oknach KPI przełączeń, detekcję dryfu pojęć w strumieniach KPI,
monitor świadomy dryfu, który ponownie uczy detektor anomalii po wykryciu
dryfu, oraz zastępczy model konfiguracja--wydajność łączący parametry HO z
wynikami mobilności.

Wyróżniają się cztery wyniki. Detektory nienadzorowane dobrze wykrywają usterki
strukturalne, z polem pod krzywą precyzja--czułość (PR-AUC) sięgającym $0,858$,
lecz w większości pomijają subtelne. Każdy
wstrzyknięty dryf zostaje ostatecznie wykryty, więc rzeczywistym wyborem
projektowym jest kompromis między opóźnieniem detekcji a fałszywymi alarmami, a
nie to, czy dryf w ogóle zostanie zauważony. Naiwne ponowne uczenie monitora na
bieżących danych przynosi efekt odwrotny: wchłania anomalie po dryfie do
``normy'' i w przybliżeniu o połowę obniża jakość detekcji, co jest efektem
samozatrucia. Prosty samonadzorowany filtr, który usuwa najbardziej podejrzane
okna przed ponownym uczeniem, naprawia to i ponad trzykrotnie zwiększa PR-AUC w
fazie dryfu. Wreszcie zastępczy model oparty na gradient boosting trafnie
przewiduje współczynnik skuteczności przełączeń (HOSR) z $R^2 = 0,96$ i,
opakowany w konformalne przedziały predykcji,
napędza kompletną pętlę: wykryj, rekomenduj, przywróć, z przewidzianym przez
model wzrostem HOSR o $+0,62$. Ogólny wniosek jest taki, że sposób aktualizowania
monitora ma większe znaczenie niż to, czy jest on w ogóle aktualizowany.

}
